\chapter*{Abstract}
\addcontentsline{toc}{chapter}{Abstract}

Physical quantities are fundamental to scientific, engineering, and medical data on the Semantic Web. The dominant approach is ontology-based: \ac{qudt} and \ac{om} provide quantity kinds, unit hierarchies, dimensional vectors, and conversion factors. But four problems persist regardless of ontological richness. The most serious is computational inaccessibility: dimensional structure is encoded in the graph but never reaches the \ac{sparql} expression evaluator, so comparisons, ordering, and aggregation over mixed-unit data produce physically incomplete results without warning. The second is bounded coverage: derived units multiply combinatorial and no enumeration can be complete. The third is interoperability: independently developed ontologies assign different \acp{iri} to the same physical unit, requiring manual cross-ontology mappings. The fourth is verbosity: a single quantity value requires at least three triples.

The \ac{lindt} specification addresses all four problems at the datatype level through \texttt{cdt:ucum}, which encodes a complete quantity as a single typed literal using \ac{ucum}. Because \ac{ucum} is a grammar rather than a list, coverage is unlimited by construction. The value space recognises dimensional equivalence natively, and operator semantics expose unit-aware equality, comparison, and arithmetic directly within \ac{sparql}. The only existing implementation is a fork of Apache Jena and it is outdated, and no implementation exists for any other major \ac{rdf} library. No general, reproducible procedure for building one has been documented.

This thesis addresses that gap. It presents a five-step methodology covering candidate library identification, extension point analysis, feasibility assessment, implementation, and validation. The feasibility step is an explicit go/no-go decision point before any code is written, and its output records not only whether each obligation can be met but through what mechanism and at what stability cost. A hard constraint frames the assessment: the implementation must leave the host library unmodified.

The methodology is validated through four working implementations across Apache Jena in Java, rdflib.js in TypeScript, dotNetRDF in C\#, and RDFLib in Python, each paired with a \ac{ucum} backend suited to its ecosystem and accompanied by a test suite written against a single shared reference checklist. The central finding is that how far an extension can reach into a host from outside determines what the implementation can deliver. That reach is set by two factors: the host library's architecture and the host language. 

The results confirm that the \ac{cdt} approach is technically viable across heterogeneous programming ecosystems and provides a lightweight, standards-compatible path to unit-aware query evaluation on the Semantic Web.